%!TEX program = lualatex
\documentclass[ngerman,headheight=70pt]{scrartcl}
\usepackage[ngerman]{babel}
\usepackage{blindtext}
\usepackage[utf8]{luainputenc}
\usepackage{scrlayer-scrpage}
\usepackage[T1]{fontenc}
\usepackage{fontspec}
\usepackage{multicol}
\usepackage{enumitem}
\usepackage[german=quotes]{csquotes}
\usepackage[bottom=2cm,footskip=8mm]{geometry}
\usepackage{pdfpages}
\setmainfont{Source Sans Pro}
\MakeOuterQuote{"}

\setlength{\parindent}{0pt}
\setlength{\parskip}{6pt}

\pagestyle{plain.scrheadings}
\thispagestyle{scrheadings}
\lohead{\vspace*{-1.5cm}\includegraphics[scale=0.4]{../images/stupa_logo_rot.pdf}}
\rohead{\\Präsidium des\\Studierendenparlaments}

\lofoot[Universität Hamburg · Präsidium des Studierendenparlaments\\\hrule Von-Melle-Park 5 · D-20146 Hamburg]{Universität Hamburg · Präsidium des Studierendenparlaments\\\hrule Von-Melle-Park 5 · D-20146 Hamburg}
\cofoot[]{}
\cohead[\thepage]{}

\renewcommand{\thesubsection}{TOP \arabic{subsection}}
\renewcommand{\thesubsubsection}{\arabic{subsubsection}.}

\begin{document}
    UHH · StuPa-Präsidium · Von-Melle-Park 5 · D-20146 Hamburg

    %%% PAGE ONE %%%%
    \section*{Protokoll der 1. Sitzung des Haushaltsausschusses vom 27. April 2016}

    \textbf{Protokoll: Jim Martens}\\
    \textbf{Ort: VMP 5 3017}\\
    \textbf{Beginn: 16.25 Uhr}\\
    \textbf{Ende: 16.38 Uhr}

    \vspace{0.5cm}
    \begin{tabular}{ll}
        Anwesend: & Melf Johannsen, Oliver Vornfeld \\
                  & Malte Jahn, Tobias Berking, Ajdina Karahasan\\
                  & \\
        Entschuldigt: & Claas-Friso Hente \\
                      &\\
        Unentschuldigt abwesend: & Nico Scharfe \\
    \end{tabular}

    %%% Tagesordnung %%%
    \underline{Tagesordnung}
    \begin{enumerate}[label={\textbf{Top \theenumi}},leftmargin=*]
        \item Begrüßung durch ein Mitglied des Präsidiums
        \item Konstituierung des Ausschusses
        \item Arbeitsplanung für die Wahlperiode 2016/2017
        \item Verschiedenes
    \end{enumerate}

    \textit{Beschlussfähigkeit mit 6 von 7 Mitgliedern festgestellt.}

    \subsection{Begrüßung durch ein Mitglied des Präsidiums}

    Ramon begrüßt die Mitglieder.

    \subsection{Konstituierung des Ausschusses}

    Melf ist einstimmig zum Vorsitzenden gewählt. Es gab
    keine weiteren Vorschläge. Abstimmung per Handzeichen.
    Melf nimmt die Wahl an.

    Oliver schlägt Tobias Berking für die Stellvertretung vor.
    Tobias Berking ist einstimmig zur Stellvertretung gewählt worden.
    Er nimmt die Wahl an.

    \subsection{Arbeitsplanung für die Wahlperiode 2016/2017}

    Tobias schlägt vor, dass Melf eine leicht verständliche Übersicht
    zur Zeitplanung anlegt. Es gibt keinen Widerspruch.

    Tobias sagt, dass der Ausschuss eine vorbereitende Aufgabe für
    das Parlament hat und dementsprechend vernünftige Pakete schnüren
    sollte. Ausschusssitzungen sollen nicht ermüdend sein.

    Melf sagt, dass sich der Ausschuss wieder trifft, wenn der/die
    Finanzreferent*In des AStA einen Haushalt vorlegt.

    \subsection{Verschiedenes}

    Es besteht kein Widerspruch gegen den Vorschlag den Verteiler zu
    öffnen.

\end{document}
